\label{sec:bot_design}

\subsection{Trading Signal and Fitness}

Each strategy produces a short-window (low-period) and a long-window
(high-period) smoothed price series. The trading signal detects crossover
events in their difference using a two-point weighted kernel:

%TC:ignore
\begin{eqnarray}
  \label{signal}
  d_t = \mathrm{Series}_{\text{short},t} - \mathrm{Series}_{\text{long},t}
  \quad\text{, }\quad
  C_t = \tfrac{1}{2}\!\left(\operatorname{sign}(d_t) -
  \operatorname{sign}(d_{t-1})\right)
\end{eqnarray}
\begin{equation}
  S =
  \left\{
  \begin{array}{ll}
    \text{buy}  & \text{if } C_t = -1 \;(\text{short crosses below long}) \\
    \text{sell} & \text{if } C_t = +1 \;(\text{short crosses above long})
  \end{array}
  \right.
\end{equation}
%TC:endignore
where $d_t$ is the difference of the two smoothed series at time $t$, $C_t
\in \{-1, 0, +1\}$ is the crossover indicator, and $S$ is the resulting
trade action. The signal is non-zero only at the moment of a crossover; it
is zero on all other timesteps.

The fitness function measures final USD holdings after simulating trades
over the full training price series:
\begin{equation} \label{eqn:fitness}
  \mathrm{Fitness}(X) = \mathrm{USD}_{\text{final}}
\end{equation}
where $X$ is the candidate parameter vector, $\mathrm{USD}_{\text{final}}$
is the cash held at the end of the simulation, and the starting balance is
$\mathrm{USD}_{\text{initial}} = \$1000$. A transaction fee of 3\% is
applied to each trade. The simulation is a single continuous pass over the
training data with no mid-run resets.

One acknowledged limitation is that the fitness function imposes no penalty
for non-trading: if no crossover ever occurs, $C_t = 0$ for all $t$ and the
bot holds cash throughout, returning $\mathrm{Fitness}(X) =
\mathrm{USD}_{\text{initial}}$. This is a valid but uninformative fixed
point. The bad/broken optimum was observed empirically --- MRFO converged
to parameter regions with no crossover signal under certain configurations.
A risk-adjusted metric such as the Sharpe ratio would penalise this
behaviour but was not adopted here; the limitation is documented rather
than silently corrected.

\subsection{Strategy Configurations}

Five strategy configurations define the hypothesis space presented to each optimiser.

\textit{Strategy 1 --- Double SMA crossover ($d=2$).} Both the short window
$x_0 \in [2,50]$ and the long window $x_1 \in [51,200]$ \cite{investopedia_moving_averages} 
are optimised. The non-overlapping constraint (short $<$ long) is guaranteed by construction
through the parameter bounds, with no runtime enforcement required. The SMA
assigns uniform weight to all prices within the window, making it
insensitive to the ordering of recent versus older prices.

\textit{Strategy 2 --- Double LMA crossover ($d=2$).} The parameter
structure is identical to Strategy 1 ($x_0 \in [2,50]$, $x_1 \in
[51,200]$). The LMA assigns linearly decaying weights so that more recent
prices contribute more than older ones, providing momentum sensitivity that
the uniform SMA cannot capture. The optimiser thus faces the same
two-dimensional search space as Strategy 1 but with a smoother series that
responds differently to regime changes.

\textit{Strategy 3 --- EMA with shared $\alpha$ ($d=3$).} A short span, a
long span, and a single shared smoothing parameter $\alpha \in (0,1)$ are
optimised jointly. Sharing one $\alpha$ value across both EMAs constrains 
both series to decay at identical rates; the
frequency contrast that drives the crossover signal therefore depends
entirely on the span difference rather than on independent decay dynamics.
This effectively collapses the degree of freedom to approximately 2.5
effective dimensions in practice.

\textit{Strategy 4 --- EMA with independent $\alpha$ ($d=4$).} Short span,
long span, $\alpha_\mathrm{short}$, and $\alpha_\mathrm{long}$ are each
optimised independently, along with additional EMA window parameters.
Decoupling the smoothing parameters gives the short-term EMA its own
reactivity and the long-term EMA its own stability, restoring the frequency
contrast that shared $\alpha$ eliminates. The expanded search space ($d=4$
versus $d=3$) carries a higher exploration cost but provides genuinely
greater expressiveness: the two EMAs can now operate at qualitatively
different timescales simultaneously. Empirically, this configuration
produced the highest Kruskal-Wallis $H$-statistics across all strategy
conditions, suggesting that the larger search space amplifies
inter-algorithm performance differences rather than masking them.

\textit{Strategy 5 --- Weighted MA combination ($d=14$).} SMA, LMA, and EMA
signals are combined as a weighted sum with learned weights $w_1$, $w_2$,
$w_3$. All three weights are treated as free optimisation dimensions and
normalised at evaluation time (each $w_i$ divided by $w_1 + w_2 + w_3$),
yielding a 14-dimensional parameter vector that includes window parameters
for each MA type plus EMA spans and smoothing factors. Deriving $w_3 = 1 -
w_1 - w_2$, as the project brief implies, would create asymmetric
optimisation pressure because $w_3$ is never directly searched; retaining
all three as free variables ensures each weight receives equal search
effort. The 14-dimensional space is the largest tested, and the curse of
dimensionality is a genuine concern when each optimiser is limited to 50 runs.

\subsection{Justified Deviations from Project Brief}

%TC:ignore
\begin{table}[H]
  \caption{\label{tab:deviations}%
    Summary of implementation decisions that deviate from the project brief,
  with justification for each.}
  \begin{ruledtabular}
    \begin{tabular}{llll}
      Deviation & \makecell[l]{Brief\\approach} &
      \makecell[l]{Our\\implementation} & Justification \\
      \hline
      EMA normalisation
      & Optional
      & \makecell[l]{Enabled by\\default}
      & \makecell[l]{Unnormalised kernels produce span-dependent\\amplitude
        distortion, making fitness values\\non-comparable across parameter
      configurations} \\[2ex]
      \makecell[l]{$d=14$ weight\\derivation}
      & \makecell[l]{$w_3$ derived as\\$1-w_1-w_2$}
      & \makecell[l]{All three weights free;\\normalised at eval time}
      & \makecell[l]{Derived $w_3$ receives no direct search
      pressure,\\creating asymmetric optimisation bias} \\[2ex]
      \makecell[l]{$n_\text{shells}$\\assignment}
      & \makecell[l]{Fixed\\shell count}
      & \makecell[l]{$\min(n_\text{shells},\,|\text{remaining}|)$}
      & \makecell[l]{Prevents degenerate empty shell states
        in\\high-dimensional runs where population may\\be smaller than the
      configured shell count} \\[2ex]
      \makecell[l]{EMA $\alpha$\\across spans}
      & \makecell[l]{Shared $\alpha$\\(Strategy~3)}
      & \makecell[l]{Independent
      $\alpha_\text{short}$,\\$\alpha_\text{long}$ (Strategy~4)}
      & \makecell[l]{Shared $\alpha$ constrains both series to
        identical\\decay rates, eliminating the frequency contrast\\that
      makes the crossover signal informative} \\
    \end{tabular}
  \end{ruledtabular}
\end{table}
%TC:endignore
